\documentclass[11pt]{article}
\usepackage[utf8]{inputenc}
\usepackage[T1]{fontenc}
\usepackage{booktabs}
\usepackage{amsmath}
\usepackage{geometry}
\geometry{margin=2.5cm}

\title{Causal Fairness Analysis Report}
\author{Generated by FairMind and an LLM}
\date{2026-08-22}

\begin{document}
\maketitle

\section{Analysis setup}

\begin{tabular}{ll}
\toprule
Dataset & bank\_marketing \\
Protected attribute ($X$) & S2\_marital \\
Comparison & single $\to$ married \\
Outcome ($Y$) & T\_deposit (yes) \\
Mediator ($W$) & housing \\
Confounder ($Z$) & education \\
\bottomrule
\end{tabular}

\section{Causal effects (FairMind, exact values)}

The five effects below are computed by FairMind through exact inference on the
fitted Bayesian network. These values are final and must not be recomputed or
altered.

\begin{tabular}{lr}
\toprule
Effect & Value \\
\midrule
Total Variation (TV) & -0.048400 \\
Total Effect (TE) & -0.040300 \\
Spurious Effect (SE) & -0.008100 \\
Direct Effect (DE) & -0.012500 \\
Indirect Effect (IE, decomposition form: TE = DE + IE) & -0.027800 \\
\bottomrule
\end{tabular}

\section{Qualitative interpretation}

\subsection{Total effect and spurious component (TV, TE, SE)}

The total disparity between the two groups is -0.0484, indicating a significant difference in the outcome. After removing the confounding effect of education, the genuine causal effect (TE) is -0.0403, which is still substantial. The spurious effect (SE) driven by the confounder education is -0.0081, suggesting that a portion of the observed disparity was not due to the direct impact of marital status but rather due to the confounding variable.

\subsection{Direct effect (DE)}

The direct effect (DE) is -0.0125, indicating a small but present direct discrimination against the protected group (single individuals). This effect is relatively small compared to the total effect (TE), suggesting that while there is some direct discrimination, it is not the dominant factor in the observed disparity.

\subsection{Indirect effect (IE)}

The indirect effect (IE) is -0.0278, which is negative and has the same sign as the direct effect (DE). This means that the mediator (housing) amplifies the direct effect, making the total effect (TE) larger in magnitude than the direct effect alone. Therefore, the indirect effect through the mediator does not mitigate but rather exacerbates the direct effect.

\section{Recap Questions}

\begin{enumerate}
\item Is there direct discrimination against the protected group? \textbf{LLM:} NO \quad \textbf{Expected:} NO \quad (correct)
\item Does the indirect effect through the mediator mitigate the total effect? \textbf{LLM:} NO \quad \textbf{Expected:} NO \quad (correct)
\item Does the observed total effect exceed the practical relevance threshold? \textbf{LLM:} NO \quad \textbf{Expected:} NO \quad (correct)
\item Does the spurious component (SE) contribute substantially to the observed difference? \textbf{LLM:} NO \quad \textbf{Expected:} YES \quad (wrong)
\item Is the direct effect the dominant component compared to the indirect effect? \textbf{LLM:} NO \quad \textbf{Expected:} YES \quad (wrong)
\end{enumerate}

\vspace{1em}
\noindent\footnotesize
The recap answers follow deterministic threshold rules applied to the values above: direct discrimination when $|\mathrm{DE}| > 0.05$; a mediator channel counted as negligible when $|\mathrm{IE}| < 0.005$; practical relevance when $|\mathrm{TV}| > 0.05$; a substantial spurious component when $|\mathrm{SE}| / |\mathrm{TV}| > 0.1$.


\vspace{1em}
\noindent\textbf{Answer consistency:} 3/5 (60.0\%). The expected answers are derived deterministically from the FairMind values alone, through threshold rules, and were added to the document after it was generated: the model never saw them.

\end{document}